\documentclass[../main.tex]{subfiles}

\begin{document}
% -----------------------------
\section{Refactoring time}
% -----------------------------
So far, we have placed all of our logic inside the main function to keep things simple
and focused on understanding the core concepts. Now, as the codebase grows, we will start
refactoring it into smaller, well-structured components to improve readability, modularity,
and maintainability.
What we are trying to achieve is following structure: 
\begin{lstlisting} [language=bash, caption={Folder structure}]
middle_end
├── analysis.py
├── constant_folding.py
├── dead_code_elimination.py
├── hc_lowering.py
├── __init__.py
├── pipeline.py
hc_dialect.py
hc_main.py
\end{lstlisting}

\begin{itemize}
    \item Create \path{middle_end} folder and given Python files
    \item Create an empty \path{__init__.py}
    \item Extract sections from \path{hc_main.py} into appropriate files
\end{itemize}

\path{pipeline.py} serves as a central integration point that orchestrates and connects
multiple compiler passes into a coherent transformation pipeline.
Inside of it, add imports:
\begin{lstlisting} [language=Python, caption={Pipeline imports}]
import inspect
from dataclasses import dataclass
from xdsl.dialects import func
from xdsl.dialects.builtin import ModuleOp
from xdsl.pattern_rewriter import (
    GreedyRewritePatternApplier,
    PatternRewriteWalker,
    RewritePattern,
)

from .analysis import analyze
from .constant_folding import FoldArithInts
from .dead_code_elimination import apply_dce
from .hc_lowering import LowerHCPattern
\end{lstlisting}
Bellow it, create configuration class that will help us to toggle specific passes easier.
\begin{lstlisting} [language=Python, caption={Configuration class}]
# -----------------------------
# Configuration object
# -----------------------------
@dataclass(frozen=True)
class MiddleEndPipelineConfig:
    run_analysis: bool = False
    debug_mode: bool = False

    apply_lowering: bool = True
    apply_constant_folding: bool = True
    apply_dce: bool = True
\end{lstlisting}
One helper function, based on previous implementation:
\begin{lstlisting} [language=Python, caption={apply\_pass helper function}]
def _apply_pass(module: ModuleOp, pattern: RewritePattern) -> None:
    applier = GreedyRewritePatternApplier([pattern()])

    # Some xdsl builds have extra kwargs; enable recursion if available.
    init_sig = inspect.signature(PatternRewriteWalker.__init__)
    kwargs = {}
    for k in ("walk_regions", "apply_recursively", "walk_into_regions"):
        if k in init_sig.parameters:
            kwargs[k] = True

    walker = PatternRewriteWalker(applier, **kwargs)

    # Walk/rewrite the function bodies explicitly
    for top_block in module.body.blocks:
        for op in top_block.ops:
            if isinstance(op, func.FuncOp):
                if hasattr(walker, "rewrite_region"):
                    walker.rewrite_region(op.body)
                elif hasattr(walker, "rewrite_op"):
                    walker.rewrite_op(op)
                else:
                    raise RuntimeError("No suitable rewrite entrypoint found on PatternRewriteWalker.")
\end{lstlisting}
And the middle end pipeline itself:
\begin{lstlisting} [language=Python, caption={Middle end pipeline class}]
# -----------------------------
# Pipeline
# -----------------------------
@dataclass
class MiddleEndPipeline:
    def __init__(self, config: MiddleEndPipelineConfig):
        self.config = config

    def apply_passes(self, module: ModuleOp) -> None:
        if self.config.apply_lowering:
            # Apply lowering
            _apply_pass(module, LowerHCPattern)
            self._print_module(module, "=== AFTER LOWERING  (arith.*) ===")

        if self.config.apply_constant_folding:
            # Apply constant folding
            _apply_pass(module, FoldArithInts)
            self._print_module(module, "=== AFTER CONSTANT FOLDING ===")


        if self.config.apply_dce:
            # Apply dead code elimination
            apply_dce(module)
            self._print_module(module, "=== AFTER DEAD CODE ELIMINATION ===")

        if self.config.run_analysis:
            # Analyze pass
            analyze(module)

    def _print_module(self, module: ModuleOp, message: str):
        if self.config.debug_mode:
            print(message + "\n")
            print(module)
\end{lstlisting}

For our main file, provide necessary imports:
\begin{lstlisting} [language=Python, caption={Middle-end imports}]
from middle_end.pipeline import (
    MiddleEndPipeline, MiddleEndPipelineConfig
)
\end{lstlisting}
Then, create wanted configuration and apply middle-end passes:
\begin{lstlisting} [language=Python, caption={Middle-end imports}]
def main():
    ctx = build_context()
    module = build_module(ctx)

    print("=== HIGH LEVEL (hc.*) ===")
    print(module)


    config = MiddleEndPipelineConfig(
        apply_lowering=True,
        apply_constant_folding=True,
        apply_dce=True,
        run_analysis=False,
        debug_mode=True,
    )

    middle_end_pipeline = MiddleEndPipeline(config)
    middle_end_pipeline.apply_passes(module)
\end{lstlisting}

MiddleEndPipeline now encapsulates the execution of middle-end compiler passes and drives them according to a configurable setup.
Its \path{apply_passes} method conditionally applies lowering, constant folding, dead code elimination,
and analysis based on flags defined in MiddleEndPipelineConfig, while \path{_print_module} provides optional
debug output after each stage to help trace how the IR evolves through the pipeline.

\end{document}
